% Part: history
% Chapter: biographies
% Section: rozsa-peter

\documentclass[../../../include/open-logic-section]{subfiles}

\begin{document}

\olfileid[pt]{his}{bio}{pet} 

\olsection{R\'ozsa P\'eter}

\olphoto{peter-rozsa}{R\'ozsa P\'eter}
 
R\'ozsa P\'eter nasceu R\'osza Politzer em Budapeste, Hungria, em 17 de
fevereiro de 1905. Ela é mais conhecida por seu trabalho sobre funções
recursivas, que foi essencial para a criação do campo da teoria da
recursão.

P\'eter cresceu em tempos políticos difíceis---a Primeira Guerra
Mundial assolou o mundo durante sua adolescência---, mas pôde frequentar
a prestigiosa Escola Feminina Maria Terezia, em Budapeste, onde se
formou em 1922. Depois, estudou na Universidade P\'azm\'any P\'eter
(mais tarde rebatizada como Universidade Lor\'and E\"otv\"os), em
Budapeste. Começou a estudar química por insistência do pai, mas depois
mudou para a matemática e se formou em 1927. Embora tivesse as
credenciais necessárias para ensinar matemática no ensino médio, a
situação econômica era então terrível, pois a Grande Depressão afetava
a economia mundial. Nesse período, P\'eter fez trabalhos ocasionais
como tutora e professora particular de matemática. Por fim, voltou à
universidade para cursar a pós-graduação em matemática. Inicialmente,
planejava trabalhar em teoria dos números, mas, depois de descobrir que
seus resultados já haviam sido demonstrados, quase abandonou a
matemática por completo. Foi incentivada a trabalhar nos teoremas da
incompletude de G\"odel e, sem saber, demonstrou vários resultados dele
de maneiras diferentes. Isso restaurou sua confiança, e P\'eter passou
a escrever seus primeiros artigos sobre teoria da recursão, inspirada
pelo programa de fundamentos de David Hilbert. Obteve o doutorado em
1935 e, em 1937, tornou-se editora do \emph{Journal of Symbolic Logic}.

Os primeiros artigos de P\'eter são amplamente reconhecidos como
contribuições fundadoras à teoria das funções recursivas. Em
\cite{Peter1935a}, ela investigou a relação entre diferentes tipos de
recursão. Em \cite{Peter1935b}, mostrou que certa função definida
recursivamente não é recursiva primitiva. Isso simplificou um resultado
anterior de Wilhelm Ackermann. A função simplificada de P\'eter é hoje
frequentemente chamada de função de Ackermann---e, às vezes, mais
apropriadamente, de função de Ackermann--P\'eter. Ela escreveu o primeiro
livro sobre a teoria das funções recursivas \citep{Peter1951}.

Apesar da importância e da influência de seu trabalho, P\'eter só
obteve um cargo docente em tempo integral em 1945. Durante a ocupação
nazista da Hungria na Segunda Guerra Mundial, P\'eter foi proibida de
ensinar por leis antissemitas. Em 1944, o governo criou um gueto judeu
em Budapeste; o gueto foi isolado do restante da cidade e vigiado por
guardas armados. P\'eter foi forçada a viver ali até sua libertação, em
1945. Depois, lecionou na Faculdade de Formação de Professores de
Budapeste e, a partir de 1955, na Universidade E\"otv\"os Lor\'and. Foi a
primeira matemática húngara a se tornar doutora acadêmica em matemática
e a primeira mulher eleita para a Academia Húngara de Ciências.

P\'{e}ter era conhecida como uma professora de matemática apaixonada,
que preferia explorar com os alunos a natureza e a beleza dos problemas
matemáticos a simplesmente dar aulas expositivas. Por isso, seus alunos
a chamavam carinhosamente de ``Tia Rosa''. P\'eter morreu em 1977, aos
~71 anos.

\begin{reading}
Para outras leituras biográficas, consulte \citep{Oconnor2014} e
\citep{Andrasfai1986}. \citet{Tamassy1994} realizou uma breve entrevista
com P\'eter. Para uma leitura divertida sobre matemática, veja o livro
de P\'eter, \emph{Playing With Infinity} \citep{Peter2010}.
\end{reading}

\end{document}
